% ============================================================
% 基礎演習Ⅱ　授業用の架空論文　共通プリアンブル
%   学術誌の体裁をまねた1段組の論文レイアウト。各論文（paper-x.tex など）の
%   冒頭で % ============================================================
% 基礎演習Ⅱ　授業用の架空論文　共通プリアンブル
%   学術誌の体裁をまねた1段組の論文レイアウト。各論文（paper-x.tex など）の
%   冒頭で % ============================================================
% 基礎演習Ⅱ　授業用の架空論文　共通プリアンブル
%   学術誌の体裁をまねた1段組の論文レイアウト。各論文（paper-x.tex など）の
%   冒頭で % ============================================================
% 基礎演習Ⅱ　授業用の架空論文　共通プリアンブル
%   学術誌の体裁をまねた1段組の論文レイアウト。各論文（paper-x.tex など）の
%   冒頭で \input{preamble} して使う。LuaLaTeX（jlreq）で組む。
% ============================================================
\documentclass[paper=a4paper,fontsize=10pt,line_length=42zw,number_of_lines=38,
  baselineskip=1.7zw,head_space=22mm,foot_space=24mm]{jlreq}

\usepackage{xcolor}
\usepackage{tabularray}
\UseTblrLibrary{booktabs}
\usepackage{caption}
\usepackage{enumitem}
\usepackage{needspace}
\usepackage{fancyhdr}
\usepackage{pgfplots}
\pgfplotsset{compat=1.18}
\usepackage[hidelinks]{hyperref}

\definecolor{ink}{HTML}{1D2330}
\definecolor{inktwo}{HTML}{4A5261}
\definecolor{inkthree}{HTML}{757D8A}
\definecolor{rulegray}{HTML}{B9BEB4}
\definecolor{papertwo}{HTML}{ECEEE8}
\definecolor{bad}{HTML}{B3261E}
\definecolor{badtint}{HTML}{F9DEDA}
\definecolor{barone}{HTML}{2747C9}
\definecolor{bartwo}{HTML}{D68A17}

% ------------------------------------------------------------
% 論文ごとに設定する情報（各論文の冒頭で \renewcommand する）
% ------------------------------------------------------------
\newcommand*\paperlabel{架空論文}         % 架空論文Ｘ など
\newcommand*\journalname{（架空の学術誌）}
\newcommand*\journalvol{}
\newcommand*\runninghead{}        % 柱に出す著者名など

% ------------------------------------------------------------
% ページスタイル
% ------------------------------------------------------------
\pagestyle{fancy}
\fancyhf{}
\renewcommand{\headrulewidth}{0.3pt}
\renewcommand{\footrulewidth}{0pt}
\fancyhead[L]{\small\color{inktwo}\journalname　\journalvol}
\fancyhead[R]{\small\color{inktwo}\runninghead}
\fancyfoot[R]{\small\thepage}
\fancyfoot[L]{\scriptsize\color{inkthree}\sffamily\paperlabel　基礎演習Ⅱ（名古屋市立大学 経済学部）の教材として作成した架空の論文。著者・調査・データ・文献はすべて架空。}
\fancypagestyle{firstpage}{%
  \fancyhf{}
  \renewcommand{\headrulewidth}{0pt}
  \fancyfoot[R]{\small\thepage}
  \fancyfoot[L]{\scriptsize\color{inkthree}\sffamily\paperlabel　基礎演習Ⅱ（名古屋市立大学 経済学部）の教材として作成した架空の論文。著者・調査・データ・文献はすべて架空。}}

% ------------------------------------------------------------
% 見出し・キャプション・箇条書き
% ------------------------------------------------------------
\ModifyHeading{section}{font=\normalsize\sffamily\bfseries,before_space=1.4\baselineskip,after_space=0.5\baselineskip,
  lines=2}
\ModifyHeading{subsection}{font=\normalsize\sffamily\bfseries,before_space=0.9\baselineskip,after_space=0.3\baselineskip,
  lines=1}
\ModifyHeading{subsubsection}{font=\normalsize\sffamily,before_space=0.6\baselineskip,after_space=0.2\baselineskip,
  lines=1}
\renewcommand{\thesection}{\arabic{section}.}
\renewcommand{\thesubsection}{\arabic{section}.\arabic{subsection}}
\renewcommand{\thesubsubsection}{\arabic{section}.\arabic{subsection}.\arabic{subsubsection}}

\DeclareCaptionLabelFormat{jp}{#1#2}
\captionsetup{labelformat=jp,labelsep=quad,font=small,labelfont={sf,bf},justification=centering,skip=4pt}
\setlist{nosep,leftmargin=1.8em}

% ------------------------------------------------------------
% 表題部　\makepapertitle{題名}{副題}{著者}{所属}
% ------------------------------------------------------------
\newcommand\makepapertitle[4]{%
  \thispagestyle{firstpage}
  \begingroup
  \noindent{\small\color{inktwo}\journalname　\journalvol}\hfill{\small\sffamily\color{inktwo}【\paperlabel】}\par
  \vspace{0.5mm}{\color{inktwo}\rule{\linewidth}{0.4pt}}\par
  \vspace{7mm}
  \begin{center}
    {\Large\sffamily\bfseries #1}\par\vspace{1.5mm}
    {\normalsize\sffamily #2}\par\vspace{6mm}
    {\normalsize #3}\par\vspace{1mm}
    {\small\color{inktwo}#4}\par
  \end{center}
  \vspace{4mm}
  \endgroup}

% 要旨とキーワード
\newenvironment{abstractbox}{%
  \begin{list}{}{\leftmargin=2\zw\rightmargin=2\zw}\item\relax\small
  {\sffamily\bfseries 要旨}\par\vspace{0.5mm}}{\end{list}\vspace{2mm}}
\newcommand\keywords[1]{{\small\hspace{2\zw}\textbf{キーワード}：#1}\par\vspace{2mm}}

% 架空であることの注記（第1ページの表題の下に置く）
\newcommand\fictionnote[1]{%
  \noindent\colorbox{badtint}{\parbox{\dimexpr\linewidth-2\fboxsep\relax}{\footnotesize\color{bad}\sffamily #1}}\par\vspace{3mm}}

% 表の共通スタイル（booktabs 風）
\NewTblrEnviron{ptab}
\SetTblrInner[ptab]{cells={font=\footnotesize},rowsep=1.2pt,colsep=5pt}
\newcommand\tnote[1]{\par\vspace{1mm}\noindent\begin{minipage}{\linewidth}\footnotesize\color{inktwo}注：#1\end{minipage}\par}

% 本文の流れの中に置く表と図（フロートにしない。頁の位置を固定するため）
%   \begin{inlinetable}{キャプション} ... \end{inlinetable}
%   \begin{inlinefigure}{キャプション} ... \end{inlinefigure}（キャプションは下）
\newenvironment{inlinetable}[1]{%
  \par\vspace{0.6\baselineskip}\noindent\begin{minipage}{\linewidth}\captionof{table}{#1}\centering\footnotesize}%
  {\end{minipage}\par\vspace{0.8\baselineskip}}
\newenvironment{inlinefigure}[1]{%
  \par\vspace{0.6\baselineskip}\noindent\begin{minipage}{\linewidth}\centering\def\figcap{#1}}%
  {\captionof{figure}{\figcap}\end{minipage}\par\vspace{0.8\baselineskip}}

% 有意水準の記号
\newcommand\sig[1]{\textsuperscript{#1}}

% 引用文献リスト
\newenvironment{references}{\begin{list}{}{\leftmargin=2\zw\itemindent=-2\zw\itemsep=1.5pt\parsep=0pt}\small}{\end{list}}
 して使う。LuaLaTeX（jlreq）で組む。
% ============================================================
\documentclass[paper=a4paper,fontsize=10pt,line_length=42zw,number_of_lines=38,
  baselineskip=1.7zw,head_space=22mm,foot_space=24mm]{jlreq}

\usepackage{xcolor}
\usepackage{tabularray}
\UseTblrLibrary{booktabs}
\usepackage{caption}
\usepackage{enumitem}
\usepackage{needspace}
\usepackage{fancyhdr}
\usepackage{pgfplots}
\pgfplotsset{compat=1.18}
\usepackage[hidelinks]{hyperref}

\definecolor{ink}{HTML}{1D2330}
\definecolor{inktwo}{HTML}{4A5261}
\definecolor{inkthree}{HTML}{757D8A}
\definecolor{rulegray}{HTML}{B9BEB4}
\definecolor{papertwo}{HTML}{ECEEE8}
\definecolor{bad}{HTML}{B3261E}
\definecolor{badtint}{HTML}{F9DEDA}
\definecolor{barone}{HTML}{2747C9}
\definecolor{bartwo}{HTML}{D68A17}

% ------------------------------------------------------------
% 論文ごとに設定する情報（各論文の冒頭で \renewcommand する）
% ------------------------------------------------------------
\newcommand*\paperlabel{架空論文}         % 架空論文Ｘ など
\newcommand*\journalname{（架空の学術誌）}
\newcommand*\journalvol{}
\newcommand*\runninghead{}        % 柱に出す著者名など

% ------------------------------------------------------------
% ページスタイル
% ------------------------------------------------------------
\pagestyle{fancy}
\fancyhf{}
\renewcommand{\headrulewidth}{0.3pt}
\renewcommand{\footrulewidth}{0pt}
\fancyhead[L]{\small\color{inktwo}\journalname　\journalvol}
\fancyhead[R]{\small\color{inktwo}\runninghead}
\fancyfoot[R]{\small\thepage}
\fancyfoot[L]{\scriptsize\color{inkthree}\sffamily\paperlabel　基礎演習Ⅱ（名古屋市立大学 経済学部）の教材として作成した架空の論文。著者・調査・データ・文献はすべて架空。}
\fancypagestyle{firstpage}{%
  \fancyhf{}
  \renewcommand{\headrulewidth}{0pt}
  \fancyfoot[R]{\small\thepage}
  \fancyfoot[L]{\scriptsize\color{inkthree}\sffamily\paperlabel　基礎演習Ⅱ（名古屋市立大学 経済学部）の教材として作成した架空の論文。著者・調査・データ・文献はすべて架空。}}

% ------------------------------------------------------------
% 見出し・キャプション・箇条書き
% ------------------------------------------------------------
\ModifyHeading{section}{font=\normalsize\sffamily\bfseries,before_space=1.4\baselineskip,after_space=0.5\baselineskip,
  lines=2}
\ModifyHeading{subsection}{font=\normalsize\sffamily\bfseries,before_space=0.9\baselineskip,after_space=0.3\baselineskip,
  lines=1}
\ModifyHeading{subsubsection}{font=\normalsize\sffamily,before_space=0.6\baselineskip,after_space=0.2\baselineskip,
  lines=1}
\renewcommand{\thesection}{\arabic{section}.}
\renewcommand{\thesubsection}{\arabic{section}.\arabic{subsection}}
\renewcommand{\thesubsubsection}{\arabic{section}.\arabic{subsection}.\arabic{subsubsection}}

\DeclareCaptionLabelFormat{jp}{#1#2}
\captionsetup{labelformat=jp,labelsep=quad,font=small,labelfont={sf,bf},justification=centering,skip=4pt}
\setlist{nosep,leftmargin=1.8em}

% ------------------------------------------------------------
% 表題部　\makepapertitle{題名}{副題}{著者}{所属}
% ------------------------------------------------------------
\newcommand\makepapertitle[4]{%
  \thispagestyle{firstpage}
  \begingroup
  \noindent{\small\color{inktwo}\journalname　\journalvol}\hfill{\small\sffamily\color{inktwo}【\paperlabel】}\par
  \vspace{0.5mm}{\color{inktwo}\rule{\linewidth}{0.4pt}}\par
  \vspace{7mm}
  \begin{center}
    {\Large\sffamily\bfseries #1}\par\vspace{1.5mm}
    {\normalsize\sffamily #2}\par\vspace{6mm}
    {\normalsize #3}\par\vspace{1mm}
    {\small\color{inktwo}#4}\par
  \end{center}
  \vspace{4mm}
  \endgroup}

% 要旨とキーワード
\newenvironment{abstractbox}{%
  \begin{list}{}{\leftmargin=2\zw\rightmargin=2\zw}\item\relax\small
  {\sffamily\bfseries 要旨}\par\vspace{0.5mm}}{\end{list}\vspace{2mm}}
\newcommand\keywords[1]{{\small\hspace{2\zw}\textbf{キーワード}：#1}\par\vspace{2mm}}

% 架空であることの注記（第1ページの表題の下に置く）
\newcommand\fictionnote[1]{%
  \noindent\colorbox{badtint}{\parbox{\dimexpr\linewidth-2\fboxsep\relax}{\footnotesize\color{bad}\sffamily #1}}\par\vspace{3mm}}

% 表の共通スタイル（booktabs 風）
\NewTblrEnviron{ptab}
\SetTblrInner[ptab]{cells={font=\footnotesize},rowsep=1.2pt,colsep=5pt}
\newcommand\tnote[1]{\par\vspace{1mm}\noindent\begin{minipage}{\linewidth}\footnotesize\color{inktwo}注：#1\end{minipage}\par}

% 本文の流れの中に置く表と図（フロートにしない。頁の位置を固定するため）
%   \begin{inlinetable}{キャプション} ... \end{inlinetable}
%   \begin{inlinefigure}{キャプション} ... \end{inlinefigure}（キャプションは下）
\newenvironment{inlinetable}[1]{%
  \par\vspace{0.6\baselineskip}\noindent\begin{minipage}{\linewidth}\captionof{table}{#1}\centering\footnotesize}%
  {\end{minipage}\par\vspace{0.8\baselineskip}}
\newenvironment{inlinefigure}[1]{%
  \par\vspace{0.6\baselineskip}\noindent\begin{minipage}{\linewidth}\centering\def\figcap{#1}}%
  {\captionof{figure}{\figcap}\end{minipage}\par\vspace{0.8\baselineskip}}

% 有意水準の記号
\newcommand\sig[1]{\textsuperscript{#1}}

% 引用文献リスト
\newenvironment{references}{\begin{list}{}{\leftmargin=2\zw\itemindent=-2\zw\itemsep=1.5pt\parsep=0pt}\small}{\end{list}}
 して使う。LuaLaTeX（jlreq）で組む。
% ============================================================
\documentclass[paper=a4paper,fontsize=10pt,line_length=42zw,number_of_lines=38,
  baselineskip=1.7zw,head_space=22mm,foot_space=24mm]{jlreq}

\usepackage{xcolor}
\usepackage{tabularray}
\UseTblrLibrary{booktabs}
\usepackage{caption}
\usepackage{enumitem}
\usepackage{needspace}
\usepackage{fancyhdr}
\usepackage{pgfplots}
\pgfplotsset{compat=1.18}
\usepackage[hidelinks]{hyperref}

\definecolor{ink}{HTML}{1D2330}
\definecolor{inktwo}{HTML}{4A5261}
\definecolor{inkthree}{HTML}{757D8A}
\definecolor{rulegray}{HTML}{B9BEB4}
\definecolor{papertwo}{HTML}{ECEEE8}
\definecolor{bad}{HTML}{B3261E}
\definecolor{badtint}{HTML}{F9DEDA}
\definecolor{barone}{HTML}{2747C9}
\definecolor{bartwo}{HTML}{D68A17}

% ------------------------------------------------------------
% 論文ごとに設定する情報（各論文の冒頭で \renewcommand する）
% ------------------------------------------------------------
\newcommand*\paperlabel{架空論文}         % 架空論文Ｘ など
\newcommand*\journalname{（架空の学術誌）}
\newcommand*\journalvol{}
\newcommand*\runninghead{}        % 柱に出す著者名など

% ------------------------------------------------------------
% ページスタイル
% ------------------------------------------------------------
\pagestyle{fancy}
\fancyhf{}
\renewcommand{\headrulewidth}{0.3pt}
\renewcommand{\footrulewidth}{0pt}
\fancyhead[L]{\small\color{inktwo}\journalname　\journalvol}
\fancyhead[R]{\small\color{inktwo}\runninghead}
\fancyfoot[R]{\small\thepage}
\fancyfoot[L]{\scriptsize\color{inkthree}\sffamily\paperlabel　基礎演習Ⅱ（名古屋市立大学 経済学部）の教材として作成した架空の論文。著者・調査・データ・文献はすべて架空。}
\fancypagestyle{firstpage}{%
  \fancyhf{}
  \renewcommand{\headrulewidth}{0pt}
  \fancyfoot[R]{\small\thepage}
  \fancyfoot[L]{\scriptsize\color{inkthree}\sffamily\paperlabel　基礎演習Ⅱ（名古屋市立大学 経済学部）の教材として作成した架空の論文。著者・調査・データ・文献はすべて架空。}}

% ------------------------------------------------------------
% 見出し・キャプション・箇条書き
% ------------------------------------------------------------
\ModifyHeading{section}{font=\normalsize\sffamily\bfseries,before_space=1.4\baselineskip,after_space=0.5\baselineskip,
  lines=2}
\ModifyHeading{subsection}{font=\normalsize\sffamily\bfseries,before_space=0.9\baselineskip,after_space=0.3\baselineskip,
  lines=1}
\ModifyHeading{subsubsection}{font=\normalsize\sffamily,before_space=0.6\baselineskip,after_space=0.2\baselineskip,
  lines=1}
\renewcommand{\thesection}{\arabic{section}.}
\renewcommand{\thesubsection}{\arabic{section}.\arabic{subsection}}
\renewcommand{\thesubsubsection}{\arabic{section}.\arabic{subsection}.\arabic{subsubsection}}

\DeclareCaptionLabelFormat{jp}{#1#2}
\captionsetup{labelformat=jp,labelsep=quad,font=small,labelfont={sf,bf},justification=centering,skip=4pt}
\setlist{nosep,leftmargin=1.8em}

% ------------------------------------------------------------
% 表題部　\makepapertitle{題名}{副題}{著者}{所属}
% ------------------------------------------------------------
\newcommand\makepapertitle[4]{%
  \thispagestyle{firstpage}
  \begingroup
  \noindent{\small\color{inktwo}\journalname　\journalvol}\hfill{\small\sffamily\color{inktwo}【\paperlabel】}\par
  \vspace{0.5mm}{\color{inktwo}\rule{\linewidth}{0.4pt}}\par
  \vspace{7mm}
  \begin{center}
    {\Large\sffamily\bfseries #1}\par\vspace{1.5mm}
    {\normalsize\sffamily #2}\par\vspace{6mm}
    {\normalsize #3}\par\vspace{1mm}
    {\small\color{inktwo}#4}\par
  \end{center}
  \vspace{4mm}
  \endgroup}

% 要旨とキーワード
\newenvironment{abstractbox}{%
  \begin{list}{}{\leftmargin=2\zw\rightmargin=2\zw}\item\relax\small
  {\sffamily\bfseries 要旨}\par\vspace{0.5mm}}{\end{list}\vspace{2mm}}
\newcommand\keywords[1]{{\small\hspace{2\zw}\textbf{キーワード}：#1}\par\vspace{2mm}}

% 架空であることの注記（第1ページの表題の下に置く）
\newcommand\fictionnote[1]{%
  \noindent\colorbox{badtint}{\parbox{\dimexpr\linewidth-2\fboxsep\relax}{\footnotesize\color{bad}\sffamily #1}}\par\vspace{3mm}}

% 表の共通スタイル（booktabs 風）
\NewTblrEnviron{ptab}
\SetTblrInner[ptab]{cells={font=\footnotesize},rowsep=1.2pt,colsep=5pt}
\newcommand\tnote[1]{\par\vspace{1mm}\noindent\begin{minipage}{\linewidth}\footnotesize\color{inktwo}注：#1\end{minipage}\par}

% 本文の流れの中に置く表と図（フロートにしない。頁の位置を固定するため）
%   \begin{inlinetable}{キャプション} ... \end{inlinetable}
%   \begin{inlinefigure}{キャプション} ... \end{inlinefigure}（キャプションは下）
\newenvironment{inlinetable}[1]{%
  \par\vspace{0.6\baselineskip}\noindent\begin{minipage}{\linewidth}\captionof{table}{#1}\centering\footnotesize}%
  {\end{minipage}\par\vspace{0.8\baselineskip}}
\newenvironment{inlinefigure}[1]{%
  \par\vspace{0.6\baselineskip}\noindent\begin{minipage}{\linewidth}\centering\def\figcap{#1}}%
  {\captionof{figure}{\figcap}\end{minipage}\par\vspace{0.8\baselineskip}}

% 有意水準の記号
\newcommand\sig[1]{\textsuperscript{#1}}

% 引用文献リスト
\newenvironment{references}{\begin{list}{}{\leftmargin=2\zw\itemindent=-2\zw\itemsep=1.5pt\parsep=0pt}\small}{\end{list}}
 して使う。LuaLaTeX（jlreq）で組む。
% ============================================================
\documentclass[paper=a4paper,fontsize=10pt,line_length=42zw,number_of_lines=38,
  baselineskip=1.7zw,head_space=22mm,foot_space=24mm]{jlreq}

\usepackage{xcolor}
\usepackage{tabularray}
\UseTblrLibrary{booktabs}
\usepackage{caption}
\usepackage{enumitem}
\usepackage{needspace}
\usepackage{fancyhdr}
\usepackage{pgfplots}
\pgfplotsset{compat=1.18}
\usepackage[hidelinks]{hyperref}

\definecolor{ink}{HTML}{1D2330}
\definecolor{inktwo}{HTML}{4A5261}
\definecolor{inkthree}{HTML}{757D8A}
\definecolor{rulegray}{HTML}{B9BEB4}
\definecolor{papertwo}{HTML}{ECEEE8}
\definecolor{bad}{HTML}{B3261E}
\definecolor{badtint}{HTML}{F9DEDA}
\definecolor{barone}{HTML}{2747C9}
\definecolor{bartwo}{HTML}{D68A17}

% ------------------------------------------------------------
% 論文ごとに設定する情報（各論文の冒頭で \renewcommand する）
% ------------------------------------------------------------
\newcommand*\paperlabel{架空論文}         % 架空論文Ｘ など
\newcommand*\journalname{（架空の学術誌）}
\newcommand*\journalvol{}
\newcommand*\runninghead{}        % 柱に出す著者名など

% ------------------------------------------------------------
% ページスタイル
% ------------------------------------------------------------
\pagestyle{fancy}
\fancyhf{}
\renewcommand{\headrulewidth}{0.3pt}
\renewcommand{\footrulewidth}{0pt}
\fancyhead[L]{\small\color{inktwo}\journalname　\journalvol}
\fancyhead[R]{\small\color{inktwo}\runninghead}
\fancyfoot[R]{\small\thepage}
\fancyfoot[L]{\scriptsize\color{inkthree}\sffamily\paperlabel　基礎演習Ⅱ（名古屋市立大学 経済学部）の教材として作成した架空の論文。著者・調査・データ・文献はすべて架空。}
\fancypagestyle{firstpage}{%
  \fancyhf{}
  \renewcommand{\headrulewidth}{0pt}
  \fancyfoot[R]{\small\thepage}
  \fancyfoot[L]{\scriptsize\color{inkthree}\sffamily\paperlabel　基礎演習Ⅱ（名古屋市立大学 経済学部）の教材として作成した架空の論文。著者・調査・データ・文献はすべて架空。}}

% ------------------------------------------------------------
% 見出し・キャプション・箇条書き
% ------------------------------------------------------------
\ModifyHeading{section}{font=\normalsize\sffamily\bfseries,before_space=1.4\baselineskip,after_space=0.5\baselineskip,
  lines=2}
\ModifyHeading{subsection}{font=\normalsize\sffamily\bfseries,before_space=0.9\baselineskip,after_space=0.3\baselineskip,
  lines=1}
\ModifyHeading{subsubsection}{font=\normalsize\sffamily,before_space=0.6\baselineskip,after_space=0.2\baselineskip,
  lines=1}
\renewcommand{\thesection}{\arabic{section}.}
\renewcommand{\thesubsection}{\arabic{section}.\arabic{subsection}}
\renewcommand{\thesubsubsection}{\arabic{section}.\arabic{subsection}.\arabic{subsubsection}}

\DeclareCaptionLabelFormat{jp}{#1#2}
\captionsetup{labelformat=jp,labelsep=quad,font=small,labelfont={sf,bf},justification=centering,skip=4pt}
\setlist{nosep,leftmargin=1.8em}

% ------------------------------------------------------------
% 表題部　\makepapertitle{題名}{副題}{著者}{所属}
% ------------------------------------------------------------
\newcommand\makepapertitle[4]{%
  \thispagestyle{firstpage}
  \begingroup
  \noindent{\small\color{inktwo}\journalname　\journalvol}\hfill{\small\sffamily\color{inktwo}【\paperlabel】}\par
  \vspace{0.5mm}{\color{inktwo}\rule{\linewidth}{0.4pt}}\par
  \vspace{7mm}
  \begin{center}
    {\Large\sffamily\bfseries #1}\par\vspace{1.5mm}
    {\normalsize\sffamily #2}\par\vspace{6mm}
    {\normalsize #3}\par\vspace{1mm}
    {\small\color{inktwo}#4}\par
  \end{center}
  \vspace{4mm}
  \endgroup}

% 要旨とキーワード
\newenvironment{abstractbox}{%
  \begin{list}{}{\leftmargin=2\zw\rightmargin=2\zw}\item\relax\small
  {\sffamily\bfseries 要旨}\par\vspace{0.5mm}}{\end{list}\vspace{2mm}}
\newcommand\keywords[1]{{\small\hspace{2\zw}\textbf{キーワード}：#1}\par\vspace{2mm}}

% 架空であることの注記（第1ページの表題の下に置く）
\newcommand\fictionnote[1]{%
  \noindent\colorbox{badtint}{\parbox{\dimexpr\linewidth-2\fboxsep\relax}{\footnotesize\color{bad}\sffamily #1}}\par\vspace{3mm}}

% 表の共通スタイル（booktabs 風）
\NewTblrEnviron{ptab}
\SetTblrInner[ptab]{cells={font=\footnotesize},rowsep=1.2pt,colsep=5pt}
\newcommand\tnote[1]{\par\vspace{1mm}\noindent\begin{minipage}{\linewidth}\footnotesize\color{inktwo}注：#1\end{minipage}\par}

% 本文の流れの中に置く表と図（フロートにしない。頁の位置を固定するため）
%   \begin{inlinetable}{キャプション} ... \end{inlinetable}
%   \begin{inlinefigure}{キャプション} ... \end{inlinefigure}（キャプションは下）
\newenvironment{inlinetable}[1]{%
  \par\vspace{0.6\baselineskip}\noindent\begin{minipage}{\linewidth}\captionof{table}{#1}\centering\footnotesize}%
  {\end{minipage}\par\vspace{0.8\baselineskip}}
\newenvironment{inlinefigure}[1]{%
  \par\vspace{0.6\baselineskip}\noindent\begin{minipage}{\linewidth}\centering\def\figcap{#1}}%
  {\captionof{figure}{\figcap}\end{minipage}\par\vspace{0.8\baselineskip}}

% 有意水準の記号
\newcommand\sig[1]{\textsuperscript{#1}}

% 引用文献リスト
\newenvironment{references}{\begin{list}{}{\leftmargin=2\zw\itemindent=-2\zw\itemsep=1.5pt\parsep=0pt}\small}{\end{list}}
